\section{Discrete Galerkin--POD Reduced-Order Model}
\label{sec:discrete_galerkin_rom}

The POD basis constructed in the previous section provides a
reduced spatial representation of the temperature-deviation
field. To obtain temperature predictions at different parameter
combinations, the governing energy equation is projected onto
this basis. Unlike direct POD reconstruction, this procedure
determines the modal coefficients from the governing equations.

The present implementation uses the discrete spatial operators
of the full-order model (FOM). Consequently, the reduced-order
model (ROM) retains the spatial discretization and boundary
treatment of the original solver. This section describes the
operator decomposition, discrete lifting, Galerkin projection,
solution of the reduced system, and comparison with the
full-order solutions.

\subsection{Decomposition of the full-order spatial operator}

The spatially discretized full-order equation can be expressed
in semi-discrete form as
\begin{equation}
    \frac{\mathrm{d}\mathbf{T}}{\mathrm{d}t}
    =
    c\mathbf{C}_h\mathbf{T}
    +\frac{1}{Pe}\mathbf{L}_h\mathbf{T}
    +\frac{Br}{Pe}\boldsymbol{\Phi}_h
    +\mathbf{b}_{\mathrm{boundary}},
    \label{eq:fom_operator_decomposition}
\end{equation}
where $\mathbf{T}$ is the vector of physical temperature
unknowns. The spatial operator is decomposed as
\begin{equation}
    \boxed{
    \mathbf{A}
    =
    c\mathbf{C}_h+\frac{1}{Pe}\mathbf{L}_h.
    }
    \label{eq:full_order_operator}
\end{equation}

Here, $\mathbf{C}_h$ and $\mathbf{L}_h$ represent the discrete
convection and diffusion operators, respectively. The
convection operator includes the negative signs associated
with the convection terms on the right-hand side of the
semi-discrete full-order equation. The diffusion operator
represents the discrete Laplacian.

Using the central-difference stencils of the full-order solver,
the convection operator is defined component-wise as
\begin{align}
    (\mathbf{C}_h\mathbf{T})_{ij}
    ={}&
    -\frac{
    u_{i+1,j}T_{i+1,j}
    -u_{i-1,j}T_{i-1,j}
    }{2\Delta x}
    \nonumber\\
    &-
    \frac{
    v_{i,j+1}T_{i,j+1}
    -v_{i,j-1}T_{i,j-1}
    }{2\Delta y}.
    \label{eq:discrete_convection_operator}
\end{align}

The discrete diffusion operator is
\begin{align}
    (\mathbf{L}_h\mathbf{T})_{ij}
    ={}&
    \frac{
    T_{i+1,j}-2T_{i,j}+T_{i-1,j}
    }{\Delta x^2}
    \nonumber\\
    &+
    \frac{
    T_{i,j+1}-2T_{i,j}+T_{i,j-1}
    }{\Delta y^2}.
    \label{eq:discrete_diffusion_operator}
\end{align}

Both operators are linear in temperature because the velocity
field is prescribed and does not depend on the temperature
solution. Their matrix representations are obtained from the
spatial stencils and index mapping used by the existing
full-order matrix assembler. The convection and diffusion
contributions are separated according to their respective
coefficients, $c$ and $1/Pe$.

The ghost-cell boundary treatment is retained when constructing
the operator matrices. Substitution of the ghost-cell relations
modifies the coefficients associated with the physical
unknowns. The affine boundary contributions are collected
separately in $\mathbf{b}_{\mathrm{boundary}}$. Thus, the
discrete operators and boundary forcing are based on the same
spatial formulation as the original full-order solver.

The viscous-dissipation contribution is represented by the
vector
\[
\boldsymbol{\Phi}_h\in\mathbb{R}^{N},
\qquad N=N_xN_y.
\]
Its components are evaluated using the discrete velocity
derivatives:
\begin{align}
    \Phi_{ij}={}&
    2\left(
    \frac{u_{i+1,j}-u_{i-1,j}}{2\Delta x}
    \right)^2
    \nonumber\\
    &+
    2\left(
    \frac{v_{i,j+1}-v_{i,j-1}}{2\Delta y}
    \right)^2
    \nonumber\\
    &+
    \left(
    \frac{u_{i,j+1}-u_{i,j-1}}{2\Delta y}
    +
    \frac{v_{i+1,j}-v_{i-1,j}}{2\Delta x}
    \right)^2.
    \label{eq:discrete_dissipation}
\end{align}
The resulting two-dimensional field is flattened using the
same indexing convention as the temperature field.

For the prescribed velocity profile
\[
u(y)=18y(1-y),\qquad v=0,
\]
the continuous dissipation function is
\[
\Phi(y)=324(1-2y)^2.
\]
Since the velocity field and computational grid are fixed,
the discrete dissipation vector is independent of the
parameters $\mu$ and $Br$.

\subsection{Discrete lifting residual and boundary contributions}

The non-homogeneous temperature boundary conditions are
handled through the reference-temperature transformation:
\begin{equation}
    \boxed{
    \mathbf{T}
    =
    \mathbf{T}_{\mathrm{ref}}+\boldsymbol{\theta}.
    }
    \label{eq:discrete_lifting}
\end{equation}
Here, $\mathbf{T}_{\mathrm{ref}}$ is the discrete representation
of $T_{\mathrm{ref}}(y)=y$, and $\boldsymbol{\theta}$ is the
temperature-deviation vector.

Substituting this decomposition into the full-order
semi-discrete equation gives
\begin{align}
    \frac{\mathrm{d}\boldsymbol{\theta}}{\mathrm{d}t}
    ={}&
    c\mathbf{C}_h\boldsymbol{\theta}
    +\frac{1}{Pe}\mathbf{L}_h\boldsymbol{\theta}
    +\frac{Br}{Pe}\boldsymbol{\Phi}_h
    +\mathbf{g}_{\mathrm{lift}},
    \label{eq:lifted_discrete_equation}
\end{align}
where the discrete lifting residual is defined as
\begin{equation}
    \boxed{
    \mathbf{g}_{\mathrm{lift}}
    =
    \mathbf{A}\mathbf{T}_{\mathrm{ref}}
    +\mathbf{b}_{\mathrm{boundary}}.
    }
    \label{eq:discrete_lifting_residual}
\end{equation}

This term accounts for the contribution of the reference
temperature when it is substituted into the discrete
full-order equation. It includes both the action of the
discrete spatial operators on the reference field and the
affine contributions introduced by the boundary closures.
In general, it must be retained unless its value has been
established for the actual discrete implementation.

The lifting residual cannot be determined solely from the
continuous differential equation. Although the continuous
reference profile $T_{\mathrm{ref}}=y$ has zero second
derivatives and no streamwise variation, the discrete
residual also depends on how the boundary conditions are
implemented. In particular, ghost-cell elimination changes
the discrete matrix coefficients and introduces affine
forcing terms. Therefore, the lifting residual must be
evaluated using the actual discrete operators and boundary
closures of the full-order solver.

Two equivalent approaches can be used to evaluate this
residual.

\paragraph{Direct evaluation using the reference field.}

In the first approach, the transport and diffusion terms
are evaluated directly using $\mathbf{T}_{\mathrm{ref}}$.
The ghost-cell values are assigned according to the original
boundary conditions before applying the discrete stencils.
For a Dirichlet boundary, the ghost-cell temperature is
calculated as
\[
T_g=2T_w-T_{\mathrm{adjacent}},
\]
where $T_w$ is the prescribed wall or inlet temperature.
For the outlet Neumann condition, the ghost-cell value is
\[
T_g=T_{\mathrm{adjacent}}.
\]
The resulting discrete operator contributions give
\begin{equation}
    \mathbf{g}_{\mathrm{lift}}^{\mathrm{direct}}
    =
    c\,\mathcal{C}_h^{\mathrm{ghost}}
    (\mathbf{T}_{\mathrm{ref}})
    +
    \frac{1}{Pe}
    \mathcal{L}_h^{\mathrm{ghost}}
    (\mathbf{T}_{\mathrm{ref}}),
    \label{eq:direct_lifting_evaluation}
\end{equation}
where the superscript indicates that the operators are
evaluated using the prescribed ghost-cell values. This
approach directly applies the discrete stencils to the
reference field and includes the boundary contributions.

\paragraph{Evaluation using the full-order matrix and boundary forcing.}

In the second approach, the residual is calculated from
the assembled spatial matrix:
\begin{equation}
    \boxed{
    \mathbf{g}_{\mathrm{lift}}^{\mathrm{matrix}}
    =
    \mathbf{A}\mathbf{T}_{\mathrm{ref}}
    +\mathbf{b}_{\mathrm{boundary}}.
    }
    \label{eq:matrix_lifting_evaluation}
\end{equation}

The affine boundary forcing is obtained separately by
applying the ghost-cell boundary conditions to a zero
physical temperature vector. Since the physical values
are zero in this calculation, a Dirichlet ghost cell
becomes $T_g=2T_w$, whereas an outlet Neumann ghost cell
remains zero. These contributions isolate the constant
terms introduced by the non-homogeneous boundary
conditions. They are collected in
$\mathbf{b}_{\mathrm{boundary}}$.

The matrix-based expression therefore combines the action
of the assembled spatial matrix on the reference field
with the affine boundary forcing. Both contributions are
necessary to reproduce the direct evaluation.

The consistency of the operator decomposition can be
examined by comparing the assembled spatial matrix with
the sum of its components:
\begin{equation}
    \boxed{
    \mathbf{A}
    \approx
    c\mathbf{C}_h+\frac{1}{Pe}\mathbf{L}_h.
    }
    \label{eq:operator_verification}
\end{equation}
Likewise, the two lifting-residual evaluations can be
compared through
\begin{equation}
    \boxed{
    \mathbf{g}_{\mathrm{lift}}^{\mathrm{direct}}
    \approx
    \mathbf{A}\mathbf{T}_{\mathrm{ref}}
    +\mathbf{b}_{\mathrm{boundary}}.
    }
    \label{eq:lifting_verification}
\end{equation}

For the present problem, the reference field is linear
in the wall-normal coordinate and independent of the
streamwise coordinate. The prescribed velocity is also
independent of $x$, with $v=0$. Consequently, the interior
convection and diffusion stencils are expected to give
zero contributions when applied to the reference field.
If the ghost-cell closures preserve this reference
profile, the complete lifting residual is also expected
to be zero. The reduced formulation nevertheless retains
the general lifting term until its contribution is
established for the discrete implementation.

\subsection{Galerkin projection onto the POD basis}

The temperature-deviation field is approximated using the
$r$ retained POD modes:
\begin{equation}
    \boxed{
    \boldsymbol{\theta}\approx\mathbf{V}_r\mathbf{a},
    }
    \label{eq:rom_approximation}
\end{equation}
where
\[
\mathbf{V}_r\in\mathbb{R}^{N\times r},
\qquad
\mathbf{a}\in\mathbb{R}^{r}.
\]
The columns of $\mathbf{V}_r$ are the retained spatial POD
modes, and $\mathbf{a}$ contains their modal coefficients.
The reconstructed temperature field is consequently
\[
\mathbf{T}\approx
\mathbf{T}_{\mathrm{ref}}+\mathbf{V}_r\mathbf{a}.
\]

Since the POD modes are obtained from the left singular
vectors of the snapshot matrix, they satisfy the
orthogonality condition
\begin{equation}
    \boxed{
    \mathbf{V}_r^T\mathbf{V}_r=\mathbf{I}_r.
    }
    \label{eq:pod_orthogonality}
\end{equation}
The present implementation uses a uniform Cartesian grid
and the ordinary Euclidean POD inner product.

Substituting the reduced approximation into
Eq.~\eqref{eq:lifted_discrete_equation} gives
\begin{align}
    \mathbf{V}_r\dot{\mathbf{a}}
    ={}&
    c\mathbf{C}_h\mathbf{V}_r\mathbf{a}
    +\frac{1}{Pe}\mathbf{L}_h\mathbf{V}_r\mathbf{a}
    \nonumber\\
    &+
    \frac{Br}{Pe}\boldsymbol{\Phi}_h
    +\mathbf{g}_{\mathrm{lift}}.
\end{align}
Multiplying from the left by $\mathbf{V}_r^T$ and using
Eq.~\eqref{eq:pod_orthogonality} yields the reduced
equation. Define the projected operators as
\begin{align}
    \mathbf{C}_r
    &=
    \mathbf{V}_r^T\mathbf{C}_h\mathbf{V}_r,
    \label{eq:reduced_convection}
    \\
    \mathbf{L}_r
    &=
    \mathbf{V}_r^T\mathbf{L}_h\mathbf{V}_r,
    \label{eq:reduced_diffusion}
    \\
    \boldsymbol{\Phi}_r
    &=
    \mathbf{V}_r^T\boldsymbol{\Phi}_h,
    \label{eq:reduced_dissipation}
    \\
    \mathbf{g}_r
    &=
    \mathbf{V}_r^T\mathbf{g}_{\mathrm{lift}}.
    \label{eq:reduced_lifting}
\end{align}

The reduced convection and diffusion matrices describe
the action of the full-order spatial operators within the
retained POD subspace. The vectors
$\boldsymbol{\Phi}_r$ and $\mathbf{g}_r$ represent the
projections of the viscous-dissipation source and the
lifting residual onto that same subspace.

The dimensions of the full-order and reduced quantities
are summarized in Table~\ref{tab:galerkin_dimensions}.

\begin{table}[htbp]
    \centering
    \caption{Dimensions of the full-order and reduced
    Galerkin operators.}
    \label{tab:galerkin_dimensions}
    \begin{tabular}{lll}
        \toprule
        \textbf{Quantity} &
        \textbf{Dimensions} &
        \textbf{Description} \\
        \midrule
        $\mathbf{C}_h$ &
        $N\times N$ &
        Full-order convection matrix \\
        $\mathbf{L}_h$ &
        $N\times N$ &
        Full-order diffusion matrix \\
        $\boldsymbol{\Phi}_h$ &
        $N\times 1$ &
        Full-order dissipation vector \\
        $\mathbf{g}_{\mathrm{lift}}$ &
        $N\times 1$ &
        Full-order lifting residual \\
        $\mathbf{V}_r$ &
        $N\times r$ &
        Retained POD basis \\
        $\mathbf{C}_r$ &
        $r\times r$ &
        Reduced convection matrix \\
        $\mathbf{L}_r$ &
        $r\times r$ &
        Reduced diffusion matrix \\
        $\boldsymbol{\Phi}_r$ &
        $r\times 1$ &
        Reduced dissipation vector \\
        $\mathbf{g}_r$ &
        $r\times 1$ &
        Reduced lifting residual \\
        \bottomrule
    \end{tabular}
\end{table}

For the present grid, $N=N_xN_y=1000$. Thus, the
full-order operators act on 1000 physical degrees of
freedom, whereas the reduced operators act only on the
$r$ retained modal coefficients.

The resulting reduced semi-discrete equation is
\begin{equation}
    \boxed{
    \frac{\mathrm{d}\mathbf{a}}{\mathrm{d}t}
    =
    c\mathbf{C}_r\mathbf{a}
    +\frac{1}{Pe}\mathbf{L}_r\mathbf{a}
    +\frac{Br}{Pe}\boldsymbol{\Phi}_r
    +\mathbf{g}_r.
    }
    \label{eq:reduced_transient_equation}
\end{equation}

\subsection{Steady reduced-order system}

The objective is to calculate steady-state temperature
fields. Accordingly, the steady-state condition is applied
to the reduced equation:
\[
\frac{\mathrm{d}\mathbf{a}}{\mathrm{d}t}=\mathbf{0}.
\]
Setting Eq.~\eqref{eq:reduced_transient_equation} to zero
and multiplying by $Pe$ gives
\begin{equation}
    cPe\mathbf{C}_r\mathbf{a}
    +\mathbf{L}_r\mathbf{a}
    +Br\boldsymbol{\Phi}_r
    +Pe\mathbf{g}_r
    =\mathbf{0}.
\end{equation}
Rearranging, the general steady reduced-order system is
\begin{equation}
    \boxed{
    (cPe\mathbf{C}_r+\mathbf{L}_r)\mathbf{a}
    =
    -Br\boldsymbol{\Phi}_r-Pe\mathbf{g}_r.
    }
    \label{eq:steady_rom_general}
\end{equation}

For the present parameter study, the convection
coefficient is fixed at $c=1$. The effective parameter is
therefore $\mu=cPe=Pe$, and the reduced system becomes
\begin{equation}
    \boxed{
    (\mu\mathbf{C}_r+\mathbf{L}_r)\mathbf{a}
    =
    -Br\boldsymbol{\Phi}_r-Pe\mathbf{g}_r.
    }
    \label{eq:steady_rom_mu}
\end{equation}

If the discrete lifting residual is negligible, its
contribution can be omitted. In that case, the reduced
system simplifies to
\begin{equation}
    \boxed{
    (\mu\mathbf{C}_r+\mathbf{L}_r)\mathbf{a}
    =
    -Br\boldsymbol{\Phi}_r.
    }
    \label{eq:steady_rom_simplified}
\end{equation}

The signs in these equations follow directly from the
semi-discrete full-order formulation. In particular,
$\mathbf{C}_h$ contains the negative convection terms
on the right-hand side of the full-order equation, while
$\mathbf{L}_h$ is the discrete Laplacian. The resulting
steady reduced system therefore has the negative
viscous-dissipation contribution on its right-hand side.

These definitions must be distinguished from the
continuous weak-form formulation introduced earlier.
There, integration by parts produces a positive
diffusion stiffness matrix, conventionally denoted by
$\mathbf{K}$, and a differently signed convection
matrix. The matrices used in the present implementation
are obtained by projecting the discrete right-hand-side
operators of the full-order solver. The continuous
weak-form matrices must therefore not be substituted
directly into Eqs.~\eqref{eq:steady_rom_general}--\eqref{eq:steady_rom_simplified}.

\subsection{Reduced system solution and temperature reconstruction}

For a specified parameter combination, the reduced
system is assembled in matrix form as
\begin{equation}
    \boxed{
    \mathbf{A}_r
    =
    cPe\mathbf{C}_r+\mathbf{L}_r,
    }
    \label{eq:rom_system_matrix}
\end{equation}
and
\begin{equation}
    \boxed{
    \mathbf{b}_r
    =
    -Br\boldsymbol{\Phi}_r-Pe\mathbf{g}_r.
    }
    \label{eq:rom_system_rhs}
\end{equation}
The modal coefficient vector is obtained by solving
\begin{equation}
    \boxed{
    \mathbf{A}_r\mathbf{a}=\mathbf{b}_r.
    }
    \label{eq:rom_linear_system}
\end{equation}
This is an $r\times r$ linear system, compared with the
$N\times N$ full-order spatial system.

After calculating the modal coefficients, the
temperature-deviation field is reconstructed using the
retained POD basis:
\[
\boldsymbol{\theta}_{\mathrm{ROM}}
=
\mathbf{V}_r\mathbf{a}.
\]
The complete reduced-order temperature field is
therefore
\begin{equation}
    \boxed{
    \mathbf{T}_{\mathrm{ROM}}
    =
    \mathbf{T}_{\mathrm{ref}}
    +\mathbf{V}_r\mathbf{a}.
    }
    \label{eq:rom_temperature_reconstruction}
\end{equation}
The reconstructed vector is then reshaped into the
two-dimensional computational field using the same
spatial indexing convention as the full-order solver.

Two numerical diagnostics are used to assess the
reduced linear solve.

First, the algebraic residual is defined as
\begin{equation}
    \boxed{
    \mathbf{r}_{\mathrm{ROM}}
    =
    \mathbf{A}_r\mathbf{a}-\mathbf{b}_r.
    }
    \label{eq:rom_algebraic_residual}
\end{equation}
Its norm measures how closely the computed modal
coefficients satisfy the reduced linear system. A small
residual indicates that the reduced equations have been
solved accurately. This is a numerical solution
diagnostic; it is not a measure of the approximation
error relative to the full-order temperature field.

Second, the condition number of the reduced system
matrix is evaluated. Using the matrix 2-norm, it is
defined as
\begin{equation}
    \boxed{
    \kappa_2(\mathbf{A}_r)
    =
    \|\mathbf{A}_r\|_2
    \|\mathbf{A}_r^{-1}\|_2
    =
    \frac{\sigma_{\max}(\mathbf{A}_r)}
    {\sigma_{\min}(\mathbf{A}_r)}.
    }
    \label{eq:rom_condition_number}
\end{equation}
The condition number indicates the sensitivity of the
reduced linear system to perturbations in its
coefficients and right-hand side. A large value
indicates a more poorly conditioned system and greater
potential sensitivity to numerical errors. It does not,
by itself, indicate whether the POD approximation is
accurate.

The algebraic residual and condition number therefore
provide information about the numerical solution of
the reduced system. The accuracy of the ROM approximation
must be assessed separately by comparing reconstructed
temperature fields against full-order solutions.

\subsection{ROM comparison with full-order solutions}

The ROM temperature field is evaluated against the
corresponding full-order solution to assess the accuracy
of the reduced representation. For each evaluated case,
the relative $L_2$ error, root-mean-square error (RMSE),
and maximum absolute error are recorded. The aggregate
results for the 100 successfully evaluated cases are
summarized in Table~\ref{tab:rom_validation_errors}.

\begin{table}[htbp]
    \centering
    \caption{Aggregate ROM errors over 100 evaluated cases.}
    \label{tab:rom_validation_errors}
    \begin{tabular}{lc}
        \toprule
        \textbf{Error metric} & \textbf{Value} \\
        \midrule
        Number of successful cases & 100 \\
        Mean relative $L_2$ error & $2.401402\times10^{-2}$ \\
        Median relative $L_2$ error & $1.705007\times10^{-2}$ \\
        Maximum relative $L_2$ error & $1.262083\times10^{-1}$ \\
        Mean RMSE & $6.498504\times10^{-2}$ \\
        Mean maximum absolute error & $1.353265\times10^{-1}$ \\
        \bottomrule
    \end{tabular}
\end{table}

The mean relative $L_2$ error across the evaluated cases
is approximately $2.40\%$, while the median is
approximately $1.71\%$. The maximum relative $L_2$ error
is approximately $12.62\%$. The difference between the
mean and median indicates that the error distribution
includes cases with substantially larger errors than
the typical case.

The cases with the smallest and largest relative
$L_2$ errors are reported in
Table~\ref{tab:rom_error_extremes}.

\begin{table}[htbp]
    \centering
    \caption{Cases with the minimum and maximum relative
    $L_2$ errors.}
    \label{tab:rom_error_extremes}
    \begin{tabular}{lcc}
        \toprule
        \textbf{Case} &
        \textbf{Relative $L_2$ error} &
        \textbf{Result} \\
        \midrule
        \texttt{train\_0001} &
        $7.040618\times10^{-6}$ &
        Minimum error \\
        \texttt{train\_0020} &
        $1.262083\times10^{-1}$ &
        Maximum error \\
        \bottomrule
    \end{tabular}
\end{table}

The results show that the ROM accuracy varies across the
evaluated parameter combinations. In particular, the
maximum-error case has a considerably larger error than
the mean and median values. Therefore, the aggregate
statistics should be considered together with the
case-wise errors when assessing the reduced model.

\subsection{Computational procedure}

The implementation follows the sequence below for a
specified parameter combination:

\begin{enumerate}
    \item Load the retained POD basis
    $\mathbf{V}_r$ and the reference temperature vector
    $\mathbf{T}_{\mathrm{ref}}$.

    \item Construct the full-order convection matrix
    $\mathbf{C}_h$, diffusion matrix $\mathbf{L}_h$,
    and viscous-dissipation vector
    $\boldsymbol{\Phi}_h$ using the spatial discretization
    and boundary treatment of the original solver.

    \item Construct the boundary forcing and evaluate
    the discrete lifting residual.

    \item Project the full-order operators and source
    terms onto the retained POD basis to obtain
    $\mathbf{C}_r$, $\mathbf{L}_r$,
    $\boldsymbol{\Phi}_r$, and $\mathbf{g}_r$.

    \item Assemble the reduced matrix
    $\mathbf{A}_r$ and right-hand-side vector
    $\mathbf{b}_r$ for the specified parameters.

    \item Solve the reduced linear system for the modal
    coefficient vector $\mathbf{a}$.

    \item Reconstruct the full temperature field using
    $\mathbf{T}_{\mathrm{ROM}}=
    \mathbf{T}_{\mathrm{ref}}+\mathbf{V}_r\mathbf{a}$.

    \item For comparison cases, evaluate the ROM error
    against the corresponding full-order solution.
\end{enumerate}

The main computational stages can be summarized as
\[
\boxed{
\begin{gathered}
\text{POD basis and reference field}\\
\downarrow\\
\text{Full-order operator components and boundary terms}\\
\downarrow\\
\text{Galerkin projection}\\
\downarrow\\
\text{Reduced system assembly}\\
\downarrow\\
\text{Solution for modal coefficients}\\
\downarrow\\
\text{Temperature reconstruction}\\
\downarrow\\
\text{Full-order comparison}
\end{gathered}
}
\]

From a computational perspective, this procedure has a
potential offline-online structure. For a fixed
computational grid and prescribed velocity field,
$\mathbf{C}_h$, $\mathbf{L}_h$, and
$\boldsymbol{\Phi}_h$ are independent of the parameter
values. Their reduced projections can therefore be
constructed once and reused for different parameter
combinations. The parameter dependence then enters
through the coefficients multiplying the reduced
operators and through the reduced right-hand side.

However, the current implementation does not yet use
a completely precomputed offline-online workflow.
The evaluation scripts reconstruct the full-order
operator components for individual parameter cases
before carrying out the reduced projection and solving
the reduced system. Thus, although the operator
decomposition identifies quantities that are reusable
for a fixed grid and velocity field, the current
implementation still repeats part of the full-order
operator construction during parameter evaluation.

\subsection{Code implementation}
\label{subsec:rom_code_implementation}

The Galerkin--POD reduced-order model is implemented in the
\texttt{F04\_galerkin\_rom} directory. The implementation is
organized into five Python scripts, each responsible for a
specific stage of the reduced-order calculation, from
full-order operator construction to parameter-case evaluation.

\begin{itemize}

    \item \texttt{f01\_operators.py}: Constructs the full-order
    operator components and source terms required for the
    Galerkin projection. Its main functions are:

    \begin{itemize}
        \item \texttt{\_assemble\_spatial\_matrix}: Builds the
        full-order spatial matrix using the existing spatial
        discretization and boundary treatment.

        \item \texttt{build\_linear\_operators}: Separates the
        convection and diffusion contributions from the
        assembled spatial matrix. The resulting operators
        satisfy
        \[
            \mathbf{A}
            =
            c\mathbf{C}_h+\frac{1}{Pe}\mathbf{L}_h.
        \]
        The diffusion operator is obtained by assembling
        the spatial matrix with $c=0$ and $Pe=1$. The
        unit-coefficient operator is
        \[
            \mathbf{A}_{\mathrm{unit}}
            =
            \mathbf{C}_h+\mathbf{L}_h,
        \]
        corresponding to $c=1$ and $Pe=1$.

        \item \texttt{build\_viscous\_dissipation\_source}:
        Calculates the parameter-independent viscous-dissipation
        vector $\boldsymbol{\Phi}_h$ from the prescribed
        velocity field and its discrete derivatives.

        \item \texttt{build\_lifting\_residual}: Calculates
        the discrete lifting residual
        \[
            \mathbf{g}_{\mathrm{lift}}
            =
            \mathbf{A}\mathbf{T}_{\mathrm{ref}}
            +\mathbf{b}_{\mathrm{boundary}}.
        \]
        This accounts for the contribution of the reference
        temperature and the affine boundary forcing.

        \item \texttt{\_assemble\_forcing\_field}: Constructs
        the forcing vector associated with the equation
        \[
            \frac{\mathrm{d}\boldsymbol{\theta}}{\mathrm{d}t}
            =
            \mathbf{A}\boldsymbol{\theta}+\mathbf{b}.
        \]
        It reuses the \texttt{build\_forcing\_field} routine
        from the direct full-order solver.

        \item \texttt{verify\_operator\_decomposition}:
        Provides checks for the operator decomposition and
        the matrix-based expression of the lifting residual.

        \item \texttt{build\_operator\_components}: Collects
        the operator matrices, viscous-dissipation source,
        and lifting contribution. It is the main interface
        for constructing the full-order components required
        by the reduced model.
    \end{itemize}

    \item \texttt{f02\_reduced\_operators.py}: Projects the
    full-order operators and source terms onto the retained
    POD basis. For a basis $\mathbf{V}_r$, the reduced
    quantities are calculated as
    \begin{align*}
        \mathbf{C}_r
        &=\mathbf{V}_r^T\mathbf{C}_h\mathbf{V}_r,\\
        \mathbf{L}_r
        &=\mathbf{V}_r^T\mathbf{L}_h\mathbf{V}_r,\\
        \boldsymbol{\Phi}_r
        &=\mathbf{V}_r^T\boldsymbol{\Phi}_h,\\
        \mathbf{g}_r
        &=\mathbf{V}_r^T\mathbf{g}_{\mathrm{lift}}.
    \end{align*}
    These reduced matrices and vectors are subsequently
    used to assemble the parameter-dependent reduced system.

    \item \texttt{f03\_solve\_reduced.py}: Handles the
    construction and solution of the reduced linear system,
    as well as temperature reconstruction. Its main functions
    are:

    \begin{itemize}
        \item \texttt{construct\_reduced\_system}: Assembles
        the reduced matrix and right-hand-side vector:
        \[
            \mathbf{A}_r
            =
            cPe\mathbf{C}_r+\mathbf{L}_r,
        \]
        \[
            \mathbf{b}_r
            =
            -Br\boldsymbol{\Phi}_r-Pe\mathbf{g}_r.
        \]
        It returns the reduced system
        $\mathbf{A}_r\mathbf{a}=\mathbf{b}_r$.

        \item \texttt{solve\_reduced\_system}: Solves the
        reduced linear system for the modal coefficient
        vector $\mathbf{a}$ and returns the solution
        together with the reduced-system residual.

        \item \texttt{reconstruct\_temperature}: Uses the
        calculated modal coefficients to reconstruct the
        temperature field:
        \[
            \boldsymbol{\theta}_{\mathrm{ROM}}
            =\mathbf{V}_r\mathbf{a},
        \]
        \[
            \mathbf{T}_{\mathrm{ROM}}
            =
            \mathbf{T}_{\mathrm{ref}}
            +\boldsymbol{\theta}_{\mathrm{ROM}}.
        \]

        \item \texttt{run\_reduced\_model}: Coordinates the
        reduced-order calculation for a specified number of
        POD modes and parameter combination. It invokes the
        reduced-system construction, solution, and
        reconstruction routines, and saves the modal
        coefficients, reduced system, and reconstructed
        temperature field.
    \end{itemize}

    \item \texttt{f04\_rom\_for\_case.py}: Provides the
    case-level evaluation of the reduced-order model. Its
    main function, \texttt{run\_rom\_for\_case}, processes
    the specified cases and retained POD rank. For each
    case, the following operations are performed:

    \begin{enumerate}
        \item The full-order temperature field is computed
        using \texttt{run\_fom\_case}, which calls
        \texttt{solve\_with\_direct}.

        \item The reduced-order solution is computed using
        \texttt{run\_rom\_case}. This involves constructing
        the full-order operator components, projecting them
        onto the POD basis, assembling and solving the
        reduced system, and reconstructing the temperature
        field.

        \item The reconstructed ROM field is compared with
        the corresponding full-order temperature field to
        calculate the error metrics.

        \item The case results and reconstructed temperature
        field are saved, and the case information is added
        to the overall results summary.
    \end{enumerate}

    After processing the cases, the script saves the
    aggregated case records in a CSV file. It also identifies
    the cases with the minimum and maximum relative
    $L_2$ errors and generates comparison plots for these
    two cases.

    \item \texttt{f05\_test\_rom.py}: Generates a parameter
    grid over the prescribed parameter domain and converts
    its parameter combinations into individual simulation
    cases. It then passes these cases to
    \texttt{run\_rom\_for\_cases} to evaluate the ROM over
    the selected parameter grid. This script provides the
    main entry point for the multi-case ROM evaluation.

\end{itemize}

The overall implementation separates the construction of
the full-order operator components, their Galerkin
projection, the solution of the reduced system, and the
evaluation of individual parameter cases. Although the
full-order operators and their reduced projections depend
only on the fixed grid and prescribed velocity field, the
current evaluation workflow reconstructs the operator
components for individual cases. Therefore, the code has
the mathematical structure needed for an offline-online
decomposition, but does not yet implement a fully
precomputed offline stage in which all reusable reduced
operators are constructed once and stored for subsequent
parameter evaluations.